\chapter[Conceptualización General]{Conceptualización General}
\label{cp:concep-general}

\insertminitoc
\parindent0pt


\section{Transporte y Movilidad Urbana}
A continuación, se exploran de forma individual cada uno de los siguientes términos.

\subsection{Transporte}
De una forma generalizada entendemos lo que es 
transporte como un proceso en el cual las 
personas se trasladan de un punto A hacia 
un punto B utilizando múltiples medios y 
redes de comunicación pudiendo ser estas 
terrestres, aéreas o marítimas. ``El \acrfull{tp} en las áreas 
urbanas ha recibido una mayor atención en 
los últimos años por su contribución a mejorar 
la sostenibilidad y la calidad de vida urbana.
El \acrshort{tp} puede resultar más atractivo al 
ofrecer una movilidad de puerta a puerta'' \cite{saif_public_2019}.
No solo se trata únicamente del momento en que algo se moviliza 
físicamente, sino de todo un régimen organizado que integra 
infraestructura (carreteras, estaciones, terminales), 
medios de transporte (automóviles, autobuses, motocicletas) y aquellos 
protagonistas sociales que los utilizan.

Con este razonamiento, se observa como el transporte urbano 
cumple con un rol importante el funcionamiento de las ciudades, 
siendo este el que facilita la conexión entre zonas residenciales, 
centro de trabajo e instituciones educativas. En donde, su eficiencia no 
se evalúa solo en la rapidez de desplazamiento, sino también en otros 
aspectos como garantizar la accesibilidad, seguridad y equidad social. 
En general, la accesibilidad se define como el acceso físico a bienes, 
servicios y destinos. En el contexto de la economía y la geografía 
urbanas, la accesibilidad se caracteriza como la facilidad para 
acceder a un área o ubicación específica \cite{saif_public_2019}.
Un sistema de transporte ineficaz o deficiente se transforma en una 
dificultad para el desarrollo económico y el bienestar de la población.

\subsection{Movilidad Urbana}
La movilidad urbana va más allá del simple transporte de un lugar a otro. 
Mientras que el transporte se centra en el movimiento físico, la movilidad urbana es un concepto 
integral que se refiere a la capacidad de los ciudadanos para acceder a bienes, servicios y oportunidades 
dentro de la ciudad, un factor crucial para la calidad de vida y el desarrollo.
Según \cite{ceder_urban_2021}, la movilidad urbana contemporánea 
se caracteriza por la congestión del tráfico, la contaminación, 
la pérdida de tiempo, el ruido y una considerable ineficiencia en 
términos de capacidad y consumo de espacio en la escala requerida 
para que una economía urbana moderna funcione productivamente.
Dicho de otra manera, La movilidad urbana sostenible es la capacidad de satisfacer 
las necesidades de la sociedad para moverse libremente, obtener un amplio 
acceso a los lugares deseados, comunicarse, negociar y establecer 
relaciones sin sacrificar otros valores \cite{costa_urban_2017}.

Desde un punto de vista útil, la movilidad posee distintas propiedades:
\begin{itemize}
    \item \textbf{Motorizada}: como el automóvil particular, la motocicleta, transporte público.
    \item \textbf{No motorizada}: incluye las caminatas, bicicletas o medios alternativos.
    \item \textbf{Multimodal}: esta se da en el momento en que se alterna entre dos o más medios
    en un mismo trayecto, 
    por ejemplo, usar bicicleta y luego transporte público.
\end{itemize}

Se dice que una ciudad tiene una buena movilidad cuando ofrece múltiples alternativas 
de desplazamiento, garantizando que la población seleccione la opción que se ajuste 
más a sus necesidades. La movilidad urbana se está volviendo autónoma, eléctrica, 
conectada y fácil de usar, teniendo en cuenta la evolución de las necesidades del consumidor 
y asumiendo los problemas de un medio ambiente limpio \cite{mavlutova_urban_2023}.

Aunque considerando lo anteriormente mencionado, esto generaría cierto choque cultural, 
ya que las personas tendrían que evaluar estas iniciativas de transporte, y las autoridades 
gobernantes tendrían que crear estos ambientes para que las opciones sean factibles.

\section{Problemáticas del tránsito vehicular en ciudades en crecimiento}
Las ciudades que pasan por un crecimiento pronunciado enfrentan cierto desafío: la falta de 
transporte aumenta a un ritmo mucho mayor a la capacidad de las infraestructuras viales. Lo que 
genera congestión, pérdida de productividad, alto consumo de combustible y contaminación.

El estado actual de la movilidad urbana tiene muchos impactos negativos en las ciudades. 
A medida que las ciudades crecen, aumenta el número de vehículos en circulación, lo que 
provoca congestión vehicular, un grave problema en muchas ciudades del mundo \cite{aprigliano_sustainable_2023}.
En Latinoamérica, dicho patrón pasa continuamente: crecimiento desordenado en la urbanización, 
incremento en el parque vehicular y un transporte público limitado o ineficiente.

Y ya que el parque vehicular ha superado las vías existentes, esta desproporción genera cuellos 
de botella en horas pico, disminución de la velocidad promedio y un aumento visible en los tiempos 
de viaje.

\section{Congestionamiento Vehicular: Definición y Causas}
El congestionamiento vehicular es una circunstancia en la que la demanda del uso de la vía supera 
la capacidad que esta puede soportar, provocando consecuentemente una reducción del flujo, aumento 
en la densidad de vehículos y trayectos más largos. Según \cite{bedoya-maya_estimating_2022}, define 
congestionamiento como la impedancia que los vehículos se imponen entre sí, debido a la 
relación velocidad-flujo, en condiciones en las que el uso de un sistema de transporte se 
aproxima a su capacidad.

Además, se han identificado nuevas dimensiones del suceso: la congestión inducida por paradas 
breves o entregas urbanas \acrfull{pudo}, los cuales generan pequeños bloqueos en 
las vías urbanas. Estos \acrshort{pudo} llegan a obstaculizar el uso de las aceras y perturbar el flujo de 
tráfico principal, lo que evidentemente generaría externalidades sociales negativas 
significativas \cite{liu_estimating_2024}.

\subsection{Causas del Congestionamiento Vehicular}
La congestión responde a una mezcla de factores, ya sean estructurales, operativos y 
conductuales, las principales causas son:

\begin{itemize}
    \item \textbf{Demanda vehicular creciente y uso intensivo del auto privado:} 
    La saturación se produce principalmente ya que el número de vehículos que 
    transitan supera la capacidad de las calles principales. En parte el aumento 
    de esta demanda, fue provocada por la \acrfull{covid-19}, en donde, \cite{Capgemini_87_2020},menciona que 
    un 81\% de consumidores de servicios de transporte, evitaba utilizarlos por 
    temas de salud y seguridad, y en cambio un 78\% optaría por utilizar sus vehículos propios.
    \item \textbf{Actividades de recoger y dejar, transporte de pasajeros:} 
    Se muestran que las breves paradas para la 
    recogida y entrega de pasajeros o la entrega de paquetes, o descarga de 
    mercancía, ocupan carriles de circulación vehicular. Ocurriendo esto de 
    manera constante en zonas donde el comercio es fuerte, ocurre cerca de 
    escuelas, estaciones de transporte, contribuyendo a la saturación urbana \cite{liu_estimating_2024}.
    \item \textbf{Diseño vial e intersecciones críticas:} La infraestructura puede 
    ser insuficiente o inexistente para drenar toda la demanda de vehículos. 
    Por cada intersección mal diseñada, semáforos mal sincronizados, o carriles 
    insuficientes se generan cuellos de botella.
    \item \textbf{Políticas de movilidad y comportamiento del conductor:} Las decisiones 
    que toman los usuarios igualmente afectan la congestión. Los factores humanos como la 
    preferencia por horarios o resistencia al cambio de elección de ruta pueden mantener la 
    congestión incluso habiendo incentivos o regulaciones. 
\end{itemize}